\subsection{Kecukupan konteks dan batas generalisasi}

Diagnosis dependency menunjukkan bahwa keputusan model dapat bergantung
pada kode yang tidak terlihat dalam diff target. Namun, hasil D88 dan D89
tidak mendukung anggapan bahwa penambahan konteks selalu memperbaiki deteksi.
Pada satu keluarga dependency, perubahan snapshot mengubah seluruh verdict;
pada kohort dependency lain, tambahan isi tidak mengubah keputusan dan
menambah penggunaan token. Temuan ini memotivasi pemeriksaan relevansi sumber,
versi yang terpasang, dan cakupan kode sebelum memperluas konteks.

Granite 3.3 8B juga tidak mengubah keputusan pada kohort D89. Model yang
lebih besar dari keluarga lain tidak otomatis mengatasi keterbatasan
informasi dalam masukan. Namun, hanya satu target MAL dan 14 kandidat HN
yang diuji, sehingga hasil ini tidak menunjukkan bahwa Granite dan Qwen
setara pada populasi pembaruan npm. Perluasan model harus disertai kohort
yang cukup beragam dan label yang dapat dipertanggungjawabkan.

Batas validitas utama adalah seleksi post-hoc D88, keterpaparan dev,
ketergantungan antartarget yang berbagi dependency, ketiadaan gold manusia,
serta cakupan HN D89 yang belum mencakup perilaku runtime sensitif.
Kedua studi menggunakan konteks kode dependency, bukan retrieval contoh KB
berlabel. Hasilnya tidak menjawab hipotesis keunggulan contoh benign sebagai
pembanding malicious, maupun efek intent terhadap FPR pada studi utama.

Evaluasi berikutnya perlu mencakup beberapa keluarga dependency dengan
provenance temporal, memisahkan versi resolved dari versi bersyarat, dan
menyertakan pembaruan sah dengan perilaku network, process, atau lifecycle
yang diverifikasi independen. Card KB harus ditulis dan diperiksa oleh dua
anotator; target, dependency terkait, dan klon yang telah terpapar perlu
dikecualikan dari KB dan test final. Audit ID bukti perlu dilengkapi penilaian
apakah kutipan mendukung alasan, karena format yang valid tidak menjamin
ketepatan penalaran.
